\documentclass[10pt,a4paper]{amsart}
\usepackage{amsmath,amssymb,amsthm,graphicx}
% Retained preamble for the Stacks Project chapter snapshot in this directory.
% The chapter's own preamble.tex is a copy of the shared Stacks Project preamble; the block
% below is sliced MECHANICALLY, line for line, from it. Nothing was added, removed or
% redefined; the wrapper class, the xr-hyper/hyperref wiring and the externaldocument list
% are omitted because the assembler supplies the document class.
%
% lines taken from preamble.tex: 3, 5, 9, 11, 12, 156, 157, 158, ... (51 lines total)
%
% package coverage restored from the chapter's own preamble (2026-10-13): xy, xy-2cell

\usepackage[all]{xy}
\xyoption{2cell}
\UseAllTwocells

\newenvironment{reference}{\comment}{\endcomment}
\newenvironment{slogan}{\comment}{\endcomment}
\newenvironment{history}{\comment}{\endcomment}
\theoremstyle{plain}
\newtheorem{theorem}[subsection]{Theorem}
\newtheorem{proposition}[subsection]{Proposition}
\newtheorem{lemma}[subsection]{Lemma}
\theoremstyle{definition}
\newtheorem{definition}[subsection]{Definition}
\newtheorem{example}[subsection]{Example}
\newtheorem{exercise}[subsection]{Exercise}
\newtheorem{situation}[subsection]{Situation}
\theoremstyle{remark}
\newtheorem{remark}[subsection]{Remark}
\newtheorem{remarks}[subsection]{Remarks}
\numberwithin{equation}{subsection}
\def\lim{\mathop{\mathrm{lim}}\nolimits}
\def\colim{\mathop{\mathrm{colim}}\nolimits}
\def\Spec{\mathop{\mathrm{Spec}}}
\def\Hom{\mathop{\mathrm{Hom}}\nolimits}
\def\Ext{\mathop{\mathrm{Ext}}\nolimits}
\def\SheafHom{\mathop{\mathcal{H}\!\mathit{om}}\nolimits}
\def\SheafExt{\mathop{\mathcal{E}\!\mathit{xt}}\nolimits}
\def\Sch{\mathit{Sch}}
\def\Mor{\mathop{\mathrm{Mor}}\nolimits}
\def\Ob{\mathop{\mathrm{Ob}}\nolimits}
\def\Sh{\mathop{\mathit{Sh}}\nolimits}
\def\NL{\mathop{N\!L}\nolimits}
\def\CH{\mathop{\mathrm{CH}}\nolimits}
\def\proetale{{pro\text{-}\acute{e}tale}}
\def\etale{{\acute{e}tale}}
\def\QCoh{\mathit{QCoh}}
\def\Ker{\mathop{\mathrm{Ker}}}
\def\Im{\mathop{\mathrm{Im}}}
\def\Coker{\mathop{\mathrm{Coker}}}
\def\Coim{\mathop{\mathrm{Coim}}}
\DeclareMathSymbol{\boxtimes}{\mathbin}{AMSa}{"02}
\def\QCohstack{\mathcal{QC}\!\mathit{oh}}
\def\Cohstack{\mathcal{C}\!\mathit{oh}}
\def\Spacesstack{\mathcal{S}\!\mathit{paces}}
\def\Quotfunctor{\mathrm{Quot}}
\def\Hilbfunctor{\mathrm{Hilb}}
\def\Curvesstack{\mathcal{C}\!\mathit{urves}}
\def\Polarizedstack{\mathcal{P}\!\mathit{olarized}}
\def\Complexesstack{\mathcal{C}\!\mathit{omplexes}}
\def\Pic{\mathop{\mathrm{Pic}}\nolimits}
\def\Picardstack{\mathcal{P}\!\mathit{ic}}
\def\Picardfunctor{\mathrm{Pic}}
\def\Deformationcategory{\mathcal{D}\!\mathit{ef}}
\usepackage{cleveref}
\crefname{originalproposition}{original proposition}{original propositions}
\Crefname{originalproposition}{Original proposition}{Original propositions}
\makeatletter
\newlabel{derived-lemma-fully-faithful-adjoint-kernel-zero}{{09J1}{source}{}{sourcepointer-1}{}}
\newlabel{derived-lemma-fully-faithful-adjoint-kernel-zero@cref}{{[originalproposition][0][]09J1}{[1][0][]source}{}{}{}}
\makeatother
\crefname{originalproposition}{original proposition}{original propositions}
\Crefname{originalproposition}{Original proposition}{Original propositions}
\makeatletter
\newlabel{lemma-QC-compare}{{lemma-QC-compare}{source}{}{sourcepointer-2}{}}
\newlabel{lemma-QC-compare@cref}{{[originalproposition][0][]lemma-QC-compare}{[1][0][]source}{}{}{}}
\makeatother
\crefname{originalproposition}{original proposition}{original propositions}
\Crefname{originalproposition}{Original proposition}{Original propositions}
\makeatletter
\newlabel{lemma-compare}{{lemma-compare}{source}{}{sourcepointer-3}{}}
\newlabel{lemma-compare@cref}{{[originalproposition][0][]lemma-compare}{[1][0][]source}{}{}{}}
\makeatother
\crefname{originalproposition}{original proposition}{original propositions}
\Crefname{originalproposition}{Original proposition}{Original propositions}
\makeatletter
\newlabel{perfect-lemma-quasi-coherence-pullback}{{08DW}{source}{}{sourcepointer-4}{}}
\newlabel{perfect-lemma-quasi-coherence-pullback@cref}{{[originalproposition][0][]08DW}{[1][0][]source}{}{}{}}
\makeatother
\crefname{originalproposition}{original proposition}{original propositions}
\Crefname{originalproposition}{Original proposition}{Original propositions}
\makeatletter
\newlabel{section-restriction-algebraic-spaces}{{section-restriction-algebraic-spaces}{source}{}{sourcepointer-5}{}}
\newlabel{section-restriction-algebraic-spaces@cref}{{[originalproposition][0][]section-restriction-algebraic-spaces}{[1][0][]source}{}{}{}}
\makeatother
\crefname{originalproposition}{original proposition}{original propositions}
\Crefname{originalproposition}{Original proposition}{Original propositions}
\makeatletter
\newlabel{sites-cohomology-lemma-cartesian-functorial}{{0GZ1}{source}{}{sourcepointer-6}{}}
\newlabel{sites-cohomology-lemma-cartesian-functorial@cref}{{[originalproposition][0][]0GZ1}{[1][0][]source}{}{}{}}
\makeatother
\crefname{originalproposition}{original proposition}{original propositions}
\Crefname{originalproposition}{Original proposition}{Original propositions}
\makeatletter
\newlabel{sites-cohomology-lemma-cartesion-plus-topology}{{0GZS}{source}{}{sourcepointer-7}{}}
\newlabel{sites-cohomology-lemma-cartesion-plus-topology@cref}{{[originalproposition][0][]0GZS}{[1][0][]source}{}{}{}}
\makeatother
\crefname{originalproposition}{original proposition}{original propositions}
\Crefname{originalproposition}{Original proposition}{Original propositions}
\makeatletter
\newlabel{spaces-perfect-lemma-DQCoh-alternative-small}{{0H06}{source}{}{sourcepointer-8}{}}
\newlabel{spaces-perfect-lemma-DQCoh-alternative-small@cref}{{[originalproposition][0][]0H06}{[1][0][]source}{}{}{}}
\makeatother
\crefname{originalproposition}{original proposition}{original propositions}
\Crefname{originalproposition}{Original proposition}{Original propositions}
\makeatletter
\newlabel{spaces-perfect-section-QC}{{0H05}{source}{}{sourcepointer-9}{}}
\newlabel{spaces-perfect-section-QC@cref}{{[originalproposition][0][]0H05}{[1][0][]source}{}{}{}}
\makeatother
\title{Stacks, Sheaves on Algebraic Stacks: proposition-QC-compare}
\author{The Stacks Project Authors}
\date{stacks-project:stacks-sheaves.tex:local-snapshot-20261012; GNU FDL 1.2 or later}
\begin{document}
\maketitle
\noindent\textit{Editorial note.} Statement span is lines 5002-5008 of sources/stacks-algebraic-stacks-sheaves/pinned-original-source.tex; proof span is lines 5010-5128, the complete author proof (117 body lines).\par\medskip
\noindent\textit{Editorial note.} Navigation stubs: 9 cited-but-not-included labels are supplied as non-load-bearing pointers with their pinned tags where the snapshot tag index has them.\par\medskip
\noindent\textit{Editorial note.} extra\_packages is empty on purpose: the retained preamble already loads xy with options and a second declaration clashes in the assembler wrapper.\par\medskip
\noindent\textit{Editorial note.} Text is retained verbatim, including the author's own typos and placeholders.\par\medskip
\noindent\textit{Editorial note.} Queue index 1 of queue\_algebraic-stacks-sheaves.json (pb=4756, stmt 5002-5008, proof 5010-5128).\par\medskip
\noindent\textit{Allowed standard prerequisites (editorial).}\begin{itemize}
\item Explicit prerequisite boundary. The cited inputs are prerequisites; the argument in the retained span is not replaced by any cited result.
\end{itemize}\medskip
\ifdefined\hypertarget\hypertarget{sourcepointer-1}{}\fi
\noindent\textit{Source reference.} Original proposition 09J1: sources/stacks-algebraic-stacks-sheaves/pinned-original-source.tex, source lines 5128--5128. Cited by the retained statement or proof but not part of the collected spans; supplied only as a navigation pointer. The argument does not depend on its content.\par\medskip
\ifdefined\hypertarget\hypertarget{sourcepointer-2}{}\fi
\noindent\textit{Source reference.} Original proposition lemma-QC-compare: sources/stacks-algebraic-stacks-sheaves/pinned-original-source.tex, source lines 4962--4968. Cited by the retained statement or proof but not part of the collected spans; supplied only as a navigation pointer. The argument does not depend on its content.\par\medskip
\ifdefined\hypertarget\hypertarget{sourcepointer-3}{}\fi
\noindent\textit{Source reference.} Original proposition lemma-compare: sources/stacks-algebraic-stacks-sheaves/pinned-original-source.tex, source lines 1430--1452. Cited by the retained statement or proof but not part of the collected spans; supplied only as a navigation pointer. The argument does not depend on its content.\par\medskip
\ifdefined\hypertarget\hypertarget{sourcepointer-4}{}\fi
\noindent\textit{Source reference.} Original proposition 08DW: sources/stacks-algebraic-stacks-sheaves/pinned-original-source.tex, source lines 5128--5128. Cited by the retained statement or proof but not part of the collected spans; supplied only as a navigation pointer. The argument does not depend on its content.\par\medskip
\ifdefined\hypertarget\hypertarget{sourcepointer-5}{}\fi
\noindent\textit{Source reference.} Original proposition section-restriction-algebraic-spaces: sources/stacks-algebraic-stacks-sheaves/pinned-original-source.tex, source lines 1421--1452. Cited by the retained statement or proof but not part of the collected spans; supplied only as a navigation pointer. The argument does not depend on its content.\par\medskip
\ifdefined\hypertarget\hypertarget{sourcepointer-6}{}\fi
\noindent\textit{Source reference.} Original proposition 0GZ1: sources/stacks-algebraic-stacks-sheaves/pinned-original-source.tex, source lines 5128--5128. Cited by the retained statement or proof but not part of the collected spans; supplied only as a navigation pointer. The argument does not depend on its content.\par\medskip
\ifdefined\hypertarget\hypertarget{sourcepointer-7}{}\fi
\noindent\textit{Source reference.} Original proposition 0GZS: sources/stacks-algebraic-stacks-sheaves/pinned-original-source.tex, source lines 5128--5128. Cited by the retained statement or proof but not part of the collected spans; supplied only as a navigation pointer. The argument does not depend on its content.\par\medskip
\ifdefined\hypertarget\hypertarget{sourcepointer-8}{}\fi
\noindent\textit{Source reference.} Original proposition 0H06: sources/stacks-algebraic-stacks-sheaves/pinned-original-source.tex, source lines 5128--5128. Cited by the retained statement or proof but not part of the collected spans; supplied only as a navigation pointer. The argument does not depend on its content.\par\medskip
\ifdefined\hypertarget\hypertarget{sourcepointer-9}{}\fi
\noindent\textit{Source reference.} Original proposition 0H05: sources/stacks-algebraic-stacks-sheaves/pinned-original-source.tex, source lines 5128--5128. Cited by the retained statement or proof but not part of the collected spans; supplied only as a navigation pointer. The argument does not depend on its content.\par\medskip
\begin{proposition}
\label{proposition-QC-compare}
Let $S$ be a scheme. Let $\mathcal{X} \to (\Sch/S)_{fppf}$ be a category
fibred in groupoids. Assume $\mathcal{X}$ is representable by an algebraic
space $X$. Then $\mathit{QC}(\mathcal{X})$ is canonically equivalent to
$D_\QCoh(\mathcal{O}_X)$.
\end{proposition}

\begin{proof}
Denote $X_{affine}$ the category of affine schemes \'etale over $X$
endowed with the chaotic topology and its structure sheaf $\mathcal{O}_X$, see
Derived Categories of Spaces, Section \ref{spaces-perfect-section-QC}.
The functor $u : X_\etale \to \mathcal{X}_\etale$ of Lemma \ref{lemma-compare}
gives rise to a functor $X_{affine} \to \mathcal{X}_{affine}$.
This is compatible with structure sheaves and produces a functor
$$
G :
\mathit{QC}(\mathcal{X}) = \mathit{QC}(\mathcal{X}_{affine}, \mathcal{O})
\longrightarrow
\mathit{QC}(X_{affine}, \mathcal{O}_X)
$$
See
Cohomology on Sites, Lemma \ref{sites-cohomology-lemma-cartesian-functorial}.
By Derived Categories of Spaces, Lemma
\ref{spaces-perfect-lemma-DQCoh-alternative-small}
the triangulated category $\mathit{QC}(X_{affine}, \mathcal{O}_X)$
is equivalent to $D_\QCoh(\mathcal{O}_X)$. Hence it suffices to
prove that $G$ is an equivalence.

\medskip\noindent
Consider the flat comparison morphisms
$\epsilon_\mathcal{X} : \mathcal{X}_{affine, \etale} \to \mathcal{X}_{affine}$
and $\epsilon_X : X_{affine, \etale} \to X_{affine}$
of ringed sites. Lemma \ref{lemma-QC-compare} and
(the proof of) Derived Categories of Spaces, Lemma
\ref{spaces-perfect-lemma-DQCoh-alternative-small}
show that the functors
$\epsilon_\mathcal{X}^*$ and $\epsilon_X^*$
identify $\mathit{QC}(\mathcal{X}_{affine}, \mathcal{O})$
and $\mathit{QC}(X_{affine}, \mathcal{O}_X)$
with subcategories
$Q_\mathcal{X} \subset D(\mathcal{X}_{affine, \etale}, \mathcal{O})$
and
$Q_X \subset D(X_{affine, \etale}, \mathcal{O}_X)$.
With these identifications
the functor $G$ in the first paragraph is induced by the functor
$$
Li_X^* = R\pi_{X, *}:
D(\mathcal{X}_{affine, \etale}, \mathcal{O})
\longrightarrow
D(X_{affine, \etale}, \mathcal{O}_X)
$$
where $i_X$ and $\pi_X$ are the morphisms from Lemma \ref{lemma-compare}
but with the \'etale sites replaced by the corresponding affine ones.
The reader can show that this replacement is permissible either
by reproving the lemma for the affine sites directly or
by using the equivalences of topoi
$\Sh(\mathcal{X}_{affine, \etale}) = \Sh(\mathcal{X}_\etale)$ and
$\Sh(X_{affine, \etale}) = \Sh(X_\etale)$.
The lemma also tells us $Li_X^*$ has a left adjoint
$$
L\pi_X^*:
D(X_{affine, \etale}, \mathcal{O}_X)
\longrightarrow
D(\mathcal{X}_{affine, \etale}, \mathcal{O})
$$
and moreover we have $Li_X^* \circ L\pi_X^* = \text{id}$
since $\pi_X \circ i_X$ is the identity. Thus it suffices
to show that (a) $L\pi_X^*$ sends $Q_X$ into $Q_\mathcal{X}$
and (b) the kernel of $Li_X^*$ is $0$.
See Derived Categories, Lemma
\ref{derived-lemma-fully-faithful-adjoint-kernel-zero}.

\medskip\noindent
Proof of (a). By Derived Categories of Spaces, Lemma
\ref{spaces-perfect-lemma-DQCoh-alternative-small}
we have $Q_X = D_\QCoh(X_{affine, \etale}, \mathcal{O}_X)$.
Let $K$ be an object of $Q_X$. Let $x$ be an object of
$\mathcal{X}_{affine, \etale}$ lying over the affine
scheme $U = p(x)$. Denote $f : U \to X$ the morphism corresponding
to $x$. Then we see that
$$
R\Gamma(x, L\pi_X^*K) = R\Gamma(U, Lf^*K)
$$
This follows from transitivity of pullbacks; see discussion in
Section \ref{section-restriction-algebraic-spaces}.
Next, suppose that $x \to x'$ is a morphism of $\mathcal{X}_{affine, \etale}$
lying over the morphism $h : U \to U'$ of affine schemes.
As before denote $f : U \to X$ and $f' : U' \to X$ the morphisms
corresponding to $x$ and $x'$ so that we have $f = f' \circ h$.
Then
\begin{align*}
R\Gamma(x, L\pi_X^*K)
& =
R\Gamma(U, Lf^*K) \\
& =
R\Gamma(U, Lh^*L(f')^*K) \\
& =
R\Gamma(U', L(f')^*K) \otimes_{\mathcal{O}(U')}^\mathbf{L} \mathcal{O}(U) \\
& =
R\Gamma(x', L\pi_X^*K) \otimes_{\mathcal{O}(x')}^\mathbf{L} \mathcal{O}(x)
\end{align*}
and hence we have (a) by the footnote in the statement of
Cohomology on Sites, Lemma
\ref{sites-cohomology-lemma-cartesion-plus-topology}. The third equality is
Derived Categories of Schemes, Lemma
\ref{perfect-lemma-quasi-coherence-pullback}.

\medskip\noindent
Proof of (b). Let $M$ be an object of $Q_\mathcal{X}$
such that $Li_X^*M = 0$. Let $x'$ be an object of $\mathcal{X}_{affine, \etale}$
lying over the affine scheme $U' = p(x')$ and assume that the corresponding
morphism $f' : U' \to X$ is \'etale. Then $f' : U' \to X$ is an object of
$X_{affine, \etale}$ and the condition $Li_X^*M = 0$ implies that
$M|_{U'_\etale} = 0$. In particular, we see that $R\Gamma(x', M) = 0$.
However, for an arbitrary object $x$ of the site
$\mathcal{X}_{affine, \etale}$ there exists a covering $\{x_i \to x\}$
such that for each $i$ there is a morphism $x_i \to x'_i$ with
$x'_i$ corresponding to an object of $X_{affine, \etale}$.
Now since $M$ is in $Q_\mathcal{X}$ we have
$$
R\Gamma(x_i, M) =
R\Gamma(x_i', M) \otimes_{\mathcal{O}(x_i')}^\mathbf{L} \mathcal{O}(x_i) =
0
$$
and we conclude that $M$ is zero as desired.
\end{proof}

\end{document}
